\section{Penalty Reformulation, Solution Algorithm, and Computational Experiments}
\label{sec:solution_algorithm}
The bilevel problem \eqref{eq:bilevel_main_a}--\eqref{eq:bilevel_main_b} cannot be solved as it stands, because the lower-level optimality condition is implicit and the upper-level feasible set is therefore defined through the solution of another optimization problem. The strategy adopted here is to remove that implicit definition by means of a value-function penalty, in the spirit of the exact penalization results of \cite{Ye1997Exact} and the recent first-order penalty algorithms of \cite{LuMei2024,Jiang2024Primal}. Concretely, recall that the lower-level value function associated with \eqref{eq:value_function_single_level} is $v(\widehat{\alpha},\widehat{\beta})=\min_{Y\in\mathcal{X}_c}\Phi_{\mathrm{c}}(Y;\widehat{\alpha},\widehat{\beta})$. Bilevel feasibility \eqref{eq:bilevel_main_b} is then equivalent to the value-function equality
\begin{equation}
\Phi_{\mathrm{c}}(X;\widehat{\alpha},\widehat{\beta})-v(\widehat{\alpha},\widehat{\beta})=0,
\label{eq:value_function_equality}
\end{equation}
because the lower-level objective $\Phi_{\mathrm{c}}$ evaluated at the design $X$ coincides with the optimal lower-level value if and only if $X$ is itself a lower-level optimum. Replacing the implicit constraint \eqref{eq:bilevel_main_b} by the explicit equality \eqref{eq:value_function_equality} is exact, but it does not yet remove the difficulty, because $v$ is itself the value function of an optimization problem and is therefore generally nonsmooth. The penalty reformulation we adopt instead promotes \eqref{eq:value_function_equality} to a soft constraint by adding it to the upper-level objective with a positive weight $\gamma>0$, which yields the single-level surrogate
\begin{equation}
\min_{(\widehat{\alpha},\widehat{\beta})\in\mathcal{H},\,X\in\mathcal{X}_c}\;
\Psi_{\gamma}(X,\widehat{\alpha},\widehat{\beta}):=
-\Pi(X)+\gamma\bigl(\Phi_{\mathrm{c}}(X;\widehat{\alpha},\widehat{\beta})-v(\widehat{\alpha},\widehat{\beta})\bigr).
\label{eq:penalized_bilevel_problem}
\end{equation}
Two facts are worth emphasizing. First, the penalized problem \eqref{eq:penalized_bilevel_problem} is a single-level problem in $(X,\widehat{\alpha},\widehat{\beta})$, whose feasible set is the simple Cartesian product $\mathcal{X}_c\times\mathcal{H}$ and whose objective combines the negative profit with a penalty on the distance to lower-level optimality. Second, for small values of $\gamma$ the penalized problem prioritizes the upper-level profit term $-\Pi(X)$, allowing $X$ to drift away from lower-level optimality if doing so increases profit; for large values of $\gamma$ the penalty dominates and pushes $X$ toward the lower-level optimal set, recovering the bilevel feasibility condition asymptotically. The classical value-function penalty results \cite{Ye1997Exact,LuMei2024} guarantee that, under mild regularity assumptions, stationary points of \eqref{eq:penalized_bilevel_problem} approach neighborhoods of bilevel-stationary points of \eqref{eq:bilevel_main_a}--\eqref{eq:bilevel_main_b} as $\gamma$ grows.

The implemented algorithm exploits this single-level penalty viewpoint while keeping the sparsity-inducing reweighting of fixed costs consistent with the current design. Two coupled iterations are needed, an inner one for the penalty and the market coefficients and an outer one for the reweighted $\ell_1$ deployment term, because the deployment term \eqref{eq:dep_term} depends on the reference capacities $(\bar S,\bar S^h)$ that should themselves track the design returned by the algorithm. Concretely, at every outer iteration the reference capacities are written to the data files of the GAMS implementation \cite{RealRojas2025AIAA}, the inner loop alternates a pure lower-level solve and a penalized solve, the market coefficients are updated by projected steps using the gap between both solves, and only at the end of the inner loop the reference capacities are refreshed with the latest lower-level solution. The pure lower-level solve corresponds to a direct evaluation of $\arg\min_{Y\in\mathcal{X}_c}\Phi_{\mathrm{c}}(Y;\widehat{\alpha},\widehat{\beta})$, which is a convex problem under fixed deployment weights and is solved through the model \texttt{cvx-ll.gms}; the penalized solve corresponds to a single-level minimization of $\Psi_{\gamma}$ at the current $(\widehat{\alpha},\widehat{\beta})$, which is implemented in \texttt{cvx-sl.gms} with the profit term added with the opposite sign and the lower-level objective coupled through $\gamma$. Both models are solved with MOSEK as the underlying conic solver. The market-coefficient update at iteration $(k,t)$ takes the form
\begin{align}
\widehat{\alpha}^{od,(k,t+1)}
&=\Pi_{[\underline{\alpha},\overline{\alpha}]}\!\left(\widehat{\alpha}^{od,(k,t)}-\mu_\alpha\,g_\alpha^{od,(k,t)}\right),
\label{eq:alpha_update}
\\
\widehat{\beta}^{od,(k,t+1)}
&=\Pi_{[\underline{\beta},\overline{\beta}]}\!\left(\widehat{\beta}^{od,(k,t)}-\mu_\beta\,g_\beta^{od,(k,t)}\right),
\label{eq:beta_update}
\end{align}
where $\Pi_{[\cdot,\cdot]}$ denotes Euclidean projection onto the admissible interval, $\mu_\alpha,\mu_\beta>0$ are step sizes, and the search directions $g_\alpha^{od},g_\beta^{od}$ are the per-market gradient-like corrections computed from the discrepancy between the lower-level and the penalized contributions of market $(o,d)$ to $\Phi_{\mathrm{c}}$. Intuitively, if a market is underemphasized by the inner lower-level model relative to the profit-oriented penalized solve, its $\widehat{\alpha}^{od}$ is reinforced; if it is overemphasized, it is reduced. The diagonal entries $(o,o)$ are reset to their nominal values because self-markets do not contribute to the network design.

The full procedure is summarized in Algorithm~\ref{alg:implementation_scheme}. The outer loop is a majorization-minimization scheme for the deployment term: at every outer iteration, the reference capacities $(\bar A,\bar S,\bar S^h)$ are used to define the convex weighted penalty in \eqref{eq:dep_term}, and at the end of the outer iteration they are refreshed with the latest lower-level solution. The inner loop is the value-function penalty algorithm: at every inner iteration, both the lower-level and the penalized models are solved, and the market coefficients are updated by projected steps. The choice of transferring the lower-level solution to the outer loop, rather than the penalized one, is deliberate, because the outer loop is meant to update the sparse approximation of the inner problem and the lower-level design is the one that approximates the value function $v(\widehat{\alpha},\widehat{\beta})$. The early-stopping criterion of the inner loop is a relative-variation test on $\Pi(X)$ that prevents unnecessary iterations once the market coefficients have stabilized.

\begin{algorithm}[t]
\caption{Penalty-based bilevel algorithm with outer sparsity updates}
\label{alg:implementation_scheme}
\KwIn{Initial market coefficients $(\widehat{\alpha},\widehat{\beta})\in\mathcal{H}$, initial reference capacities $\bar A,\bar S,\bar S^h$, step sizes $\mu_\alpha,\mu_\beta$, penalty parameter $\gamma$, numbers of outer and inner iterations $K_{\mathrm{out}},K_{\mathrm{in}}$, tolerance $\eta>0$.}
\For{$k=1,\dots,K_{\mathrm{out}}$}{
  build the deployment term $\mathrm{DEP}$ in \eqref{eq:dep_term} from the current references $\bar A,\bar S,\bar S^h$\;
  \For{$t=1,\dots,K_{\mathrm{in}}$}{
    solve the lower-level model $\min_{Y\in\mathcal{X}_c}\Phi_{\mathrm{c}}(Y;\widehat{\alpha},\widehat{\beta})$ via \texttt{cvx-ll.gms} and recover $X^{\mathrm{ll}}$\;
    evaluate the per-market lower-level contributions of $\Phi_{\mathrm{c}}$ at $X^{\mathrm{ll}}$\;
    solve the penalized model $\min \Psi_\gamma(X,\widehat{\alpha},\widehat{\beta})$ via \texttt{cvx-sl.gms} and recover $X^{\mathrm{pen}}$\;
    compute the projected update directions $g_\alpha^{od},g_\beta^{od}$ from the difference between the lower-level and penalized market contributions\;
    update $(\widehat{\alpha},\widehat{\beta})$ via \eqref{eq:alpha_update}--\eqref{eq:beta_update}\;
    \If{$|\Pi(X^{\mathrm{pen},(t)})-\Pi(X^{\mathrm{pen},(t-1)})|/|\Pi(X^{\mathrm{pen},(t-1)})|<\eta$}{break\;}
  }
  set $\bar A,\bar S,\bar S^h\leftarrow A^{\mathrm{ll}},S^{\mathrm{ll}},S^{h,\mathrm{ll}}$\;
}
\KwOut{Sparse capacities $S^\star,S^{h,\star},A^\star$, market shares $F^\star,F^{\mathrm{ext},\star}$, routing pattern $\{F_{ij}^{od,\star}\}$, and calibrated market coefficients $(\widehat{\alpha}^\star,\widehat{\beta}^\star)$.}
\end{algorithm}

The whole scheme should be read as a nested successive-approximation algorithm. The outer MM loop progressively improves the convex surrogate of the $\ell_0$ activation cost by reweighting the linear penalties on $S_i,S_i^h$ around the latest design, and the inner penalty loop progressively improves the alignment between the lower-level model and the airline profit objective by recalibrating the market coefficients $(\widehat{\alpha},\widehat{\beta})$. Their combined action is what makes the approach computationally viable on instances that are well beyond the reach of the mixed-integer baseline.

\paragraph{Computational experiments.}
\label{sec:numerical_experiments}
The methodology is validated on four progressively larger Spanish-airline instances and on a realistic Madrid-centred international instance with Iberia-style coverage. The smallest instances are designed to allow a head-to-head comparison against a state-of-the-art commercial branch-and-bound solver, whereas the largest one is designed to illustrate the scalability of the proposed scheme on a setting where the mixed-integer baseline is no longer tractable.

The first instance is a 4-node Spanish network (\texttt{4node\_spain/}) with airports MAD, BCN, PMI and AGP and twelve directed flight legs, used as a sanity check on a setting in which the MIP can still be solved to optimality within minutes. The second instance is a 6-node Spanish network (\texttt{6node\_spain/}) that adds ALC and LPA to the previous airport set, generating thirty directed flight legs and roughly three dozen origin-destination markets. The third instance is an 8-node Spanish network (\texttt{8node\_spain/}) that further adds TFS and IBZ, producing fifty-six directed legs and substantially denser interaction patterns. All Spanish instances share the same demand source, namely the public matrix CO261 of inter-airport passenger flows in Spain, and use the official YieldPKT data of the Spanish Ministry of Transport for fares; the underlying distance matrix is reconstructed from the official airport coordinates. Finally, the largest instance is the Madrid-centred international network (\texttt{red\_iberia\_MAD/}), which contains twenty-five candidate airports including the major Iberia destinations in Europe, North and South America (MAD, BCN, AGP, AMS, BIO, BRU, EZE, FCO, FRA, GRU, IBZ, JFK, LHR, LIM, LIS, LPA, MAH, MEX, MUC, MXP, ORY, PMI, SVQ, TFN, VLC) and the corresponding international demand and fare data.

In all four instances the proposed bilevel-sparsity scheme is implemented in Python coupled with the GAMS models \texttt{cvx-ll.gms} and \texttt{cvx-sl.gms}, and the resulting convex subproblems are solved with MOSEK. The mixed-integer baseline is the model of Section~\ref{sec:problem_setting} solved with a commercial MILP solver under the same input data and a time limit of two hours per instance. For every Spanish instance, we report the mean wall-clock solution time, the number of activated airports and hubs, and the relative objective gap $|\Pi^{\mathrm{BL}}-\Pi^{\mathrm{MIP}}|/|\Pi^{\mathrm{MIP}}|$ obtained when both designs are evaluated under the original mixed-integer profit \eqref{eq:profit_mip_view}; for the international Iberia instance, only the bilevel scheme produces solutions, since the mixed-integer baseline does not return a feasible incumbent within the time limit even for the smallest budget regime considered. The aggregated comparison files \texttt{compare\_bilevel\_vs\_mip\_6node.csv} (and the corresponding 4-node and 8-node versions) collect these metrics for a sweep over budgets, hub-cost weights $\lambda$, market-coefficient initializations and step-size choices.

The qualitative findings are consistent across instance sizes. On the 4-node instance, the bilevel scheme matches the mixed-integer optimum on most budget regimes and runs in fractions of a second once the GAMS subproblems are warmed up, whereas the MIP requires tens of seconds to a few minutes. On the 6-node and 8-node instances the bilevel scheme remains in the second-to-minute regime, while the MIP frequently approaches the two-hour time limit; the resulting objective gaps depend strongly on the budget and on the choice of market-coefficient step sizes, with several configurations achieving small relative gaps and others showing the typical penalty-method behavior of trading optimality for runtime. The hub structure recovered by the bilevel scheme is qualitatively coherent with the one selected by the MIP, with MAD acting as the dominant hub in all instances and with secondary hubs appearing only when the budget is sufficient to support them. On the international Iberia instance, the bilevel scheme returns sparse network designs with a small number of European hubs (typically MAD and one of LHR, FRA or AMS) and a coherent allocation of long-haul services to the Americas; representative network maps are stored as PNG outputs in \texttt{red\_iberia\_MAD/plots\_blo/}.

The overall picture that emerges from the experiments is twofold. On the one hand, the proposed continuous bilevel scheme is several orders of magnitude faster than the mixed-integer baseline on the small Spanish instances, and is the only one of the two that produces network designs at all on the realistic Madrid-centred international instance. On the other hand, the resulting designs depend nontrivially on the bilevel hyperparameters, and a careful calibration of the step sizes $\mu_\alpha,\mu_\beta$ and of the penalty weight $\gamma$ is needed to keep the airline profit close to the mixed-integer benchmark. Both observations are in line with the behavior of value-function penalty methods reported in the recent first-order bilevel optimization literature \cite{LuMei2024,Jiang2024Primal} and motivate the lines of further work discussed in the conclusions.
